\documentclass[11pt]{article}

\usepackage[margin=1in]{geometry}
\usepackage[T1]{fontenc}
\usepackage[utf8]{inputenc}
\usepackage{lmodern}
\usepackage{booktabs}
\usepackage{longtable}
\usepackage{array}
\usepackage{xcolor}
\usepackage{hyperref}
\usepackage{enumitem}
\usepackage{titlesec}

\hypersetup{
  colorlinks=true,
  linkcolor=blue!50!black,
  urlcolor=blue!50!black
}

\setlist[itemize]{leftmargin=*, itemsep=2pt, topsep=4pt}
\setlist[enumerate]{leftmargin=*, itemsep=2pt, topsep=4pt}
\titleformat{\section}{\Large\bfseries}{\thesection}{0.6em}{}
\titleformat{\subsection}{\large\bfseries}{\thesubsection}{0.6em}{}
\renewcommand{\arraystretch}{1.18}

\newcolumntype{L}[1]{>{\raggedright\arraybackslash}p{#1}}
\newcommand{\code}[1]{\nolinkurl{#1}}

\title{\textbf{SPR 4862 Implementation Plan and Initial Data Assessment}\\
\large INDOT Solar Development Suitability Map}
\author{Md Kamruzzaman Kamrul}
\date{May 18, 2026}

\begin{document}
\maketitle

\begin{abstract}
This document consolidates the implementation plan for the INDOT Solar Development
Suitability Map with the first technical assessment performed on the local shapefile
deliverables. It records the intended system architecture, phased work plan, initial
data inventory, what worked well, what did not work as expected, and the likely
reasons behind the observed data issues. The purpose is to create a traceable baseline
before moving into data cleaning, normalization, API development, and public map
implementation.
\end{abstract}

\tableofcontents
\newpage

\section{Project Purpose}

The project goal is to build a publicly accessible, interactive web-based map for
potential solar development sites on INDOT-owned facilities and right-of-way parcels.
The application should allow users to inspect candidate sites, compare suitability
attributes, filter or search by relevant criteria, and export or review site data in
a maintainable format.

The immediate implementation goal is to turn the provided GIS shapefiles into a clean,
web-ready data layer. Later implementation can then add an API, map interface,
configuration system, deployment workflow, documentation, and handoff package.

\section{Source Materials Reviewed}

\begin{longtable}{L{0.32\textwidth} L{0.28\textwidth} L{0.32\textwidth}}
\toprule
\textbf{Source} & \textbf{Location} & \textbf{Use in Assessment} \\
\midrule
\endhead
Implementation plan & \code{Implementation_Plan_SPR4862.md} & Baseline scope, phases, functional requirements, success criteria, and deliverables. \\
Candidate site shapefile & \code{All_Candidate_Sites/All_Candidate_Sites_Final.shp} & Base candidate site inventory. \\
Facility scored shapefile & \code{Facility_Scored/solar_potential_scored_indotfacility.shp} & Scored INDOT facility subset. \\
ROW scored shapefile & \code{ROW_Scored/solar_potential_scored_interchange.shp} & Scored interchange/right-of-way subset. \\
Data inventory report & \code{docs/data_inventory.md} & Generated field-level inventory from the local shapefiles. \\
Reference report PDF & \code{R1.IDOT_dot_57410_DS1 (1).pdf} & Reference decision-support tool and suitability-scoring context. \\
Conversation notes & \code{conversation_with_Dr.Chang.txt} & Project background and Illinois tool reference. \\
\bottomrule
\end{longtable}

\section{Full Implementation Plan}

\subsection{System Architecture}

The planned system has five major layers:

\begin{enumerate}
  \item \textbf{Data processing pipeline}: read shapefiles, validate geometry, normalize fields, repair known issues, and export clean spatial data.
  \item \textbf{Spatial data store}: start with file-based GeoJSON/CSV for the prototype; move to PostGIS for production querying if needed.
  \item \textbf{Backend API}: provide endpoints for sites, site details, statistics, configuration, and export.
  \item \textbf{Frontend map}: interactive Leaflet/React web application with layers, popups, filters, legends, and responsive behavior.
  \item \textbf{Deployment and documentation}: Dockerized environment, repeatable build/deploy steps, user guide, admin guide, and technical report.
\end{enumerate}

\subsection{Phased Delivery Plan}

\begin{longtable}{L{0.18\textwidth} L{0.28\textwidth} L{0.44\textwidth}}
\toprule
\textbf{Phase} & \textbf{Focus} & \textbf{Primary Outputs} \\
\midrule
\endhead
1 & Foundation and data preparation & Repository structure, dependency file, shapefile inventory, schema report, data-quality findings. \\
2 & Data cleaning and export & Valid geometry, derived latitude/longitude, normalized display fields, clean GeoJSON/CSV outputs. \\
3 & Backend API & Site-list endpoint, site-detail endpoint, summary-statistics endpoint, filtering and export support. \\
4 & Frontend map prototype & Interactive base map, facility and ROW layers, popups, filters, legend, responsive layout. \\
5 & Integration and QA & End-to-end workflow tests, data validation checks, performance checks, browser review. \\
6 & Deployment and documentation & Deployment configuration, public build, user/admin docs, final technical handoff. \\
\bottomrule
\end{longtable}

\subsection{Functional Requirements}

\begin{itemize}
  \item Parse and validate all three shapefile inputs.
  \item Preserve raw source fields while adding cleaned public-facing fields.
  \item Convert projected geometries from \code{EPSG:26916} to \code{EPSG:4326} for web mapping.
  \item Display candidate sites on an interactive map.
  \item Support layer toggling for facility and ROW/interchange sites.
  \item Support filtering by site type, score fields, land area, district/layer, flood status, and related attributes.
  \item Show site details in map popups or side panels.
  \item Export filtered data as GeoJSON or CSV.
  \item Provide documentation for setup, data updates, and maintenance.
\end{itemize}

\subsection{Non-Functional Requirements}

\begin{itemize}
  \item Keep the data processing reproducible from source shapefiles.
  \item Keep the frontend responsive and usable on desktop and mobile screens.
  \item Keep map rendering fast enough for the current 100-site scale, with room for larger future data.
  \item Use a maintainable configuration approach so labels, filters, and display fields can change without deep code edits.
  \item Make the deployment path simple enough for future handoff.
\end{itemize}

\subsection{Recommended Next Implementation Order}

\begin{enumerate}
  \item Data cleaning and normalized export.
  \item Minimal static map prototype consuming the cleaned GeoJSON.
  \item Backend API only after the cleaned data model is stable.
  \item React/Leaflet interface with filters and popups.
  \item Production packaging, testing, and documentation.
\end{enumerate}

\section{Initial Assessment Performed}

\subsection{Environment Setup}

The workspace already contained a Python virtual environment at \code{.venv}, but it
did not contain GIS libraries. A minimal \code{requirements.txt} was added and installed
successfully. The dependencies installed cleanly and allowed local reading of all
three shapefiles using GeoPandas.

\begin{longtable}{L{0.32\textwidth} L{0.58\textwidth}}
\toprule
\textbf{Item} & \textbf{Result} \\
\midrule
\endhead
Python version & Python 3.10.11 \\
GIS dependencies & GeoPandas, Pyogrio, Shapely, PyProj, Pandas \\
Inventory script & \code{scripts/inspect_shapefiles.py} \\
Generated inventory & \code{docs/data_inventory.md} \\
\bottomrule
\end{longtable}

\subsection{Dataset Inventory}

\begin{longtable}{L{0.27\textwidth} L{0.13\textwidth} L{0.20\textwidth} L{0.27\textwidth}}
\toprule
\textbf{Dataset} & \textbf{Records} & \textbf{CRS} & \textbf{Geometry Types} \\
\midrule
\endhead
\code{All_Candidate_Sites_Final.shp} & 104 & \code{EPSG:26916} & Polygon, MultiPolygon \\
\code{solar_potential_scored_indotfacility.shp} & 58 & \code{EPSG:26916} & Polygon, MultiPolygon \\
\code{solar_potential_scored_interchange.shp} & 45 & \code{EPSG:26916} & Polygon, MultiPolygon \\
\bottomrule
\end{longtable}

The scored facility and ROW files together contain 103 records. The all-candidate
file contains 104 records. All files use the same projected coordinate system,
which is suitable for GIS operations but must be transformed for browser mapping.

\section{What Worked Well}

\begin{longtable}{L{0.28\textwidth} L{0.60\textwidth}}
\toprule
\textbf{Area} & \textbf{Assessment} \\
\midrule
\endhead
File availability & All three shapefiles named in the plan are present locally, including supporting \code{.dbf}, \code{.prj}, \code{.shx}, and metadata files. \\
Coordinate reference system & All datasets consistently use \code{EPSG:26916}, so conversion to web coordinates can be handled consistently. \\
Primary identifiers & \code{SPR_ID} is present in all three datasets and is unique within each dataset. This gives us a usable join/reference key. \\
Geometry coverage & The source geometries are polygonal, which is appropriate for parcel/site display and area-based map interaction. \\
Score fields & The scored facility and ROW files include suitability component fields such as \code{sol_s}, \code{slp_s}, \code{trn_s}, \code{evp_s}, \code{dem_s}, \code{fld_s}, and \code{lc_s}. \\
Metadata lineage & The ArcGIS metadata includes processing history, export dates, source feature-class names, and field-mapping information. This helped explain several unexpected data behaviors. \\
Reference context & The Illinois decision-support report explains the intended style of tool and why different property types may use different suitability logic. \\
\bottomrule
\end{longtable}

\section{What Did Not Work or Was Unexpected}

\subsection{Latitude and Longitude Fields Are Not Usable}

The \code{Lat} and \code{Long} fields are present in every dataset, but all checked
values are \code{0.0}. This means they cannot be used for map placement, popups, or
exports.

\textbf{Likely reason:} The metadata shows these fields were carried through as
ordinary attribute fields during ArcGIS export. They appear to have been placeholders
or stale fields that were not recalculated after the final geometry processing.

\textbf{Implementation response:} Ignore the raw \code{Lat} and \code{Long} fields for
location. Derive centroid longitude/latitude from geometry after transforming to
\code{EPSG:4326}. Preserve the raw fields only for traceability if needed.

\subsection{One Candidate Site Has No Scored Counterpart}

The all-candidate dataset contains 104 records, while the two scored layers contain
103 total records. The missing scored record is:

\begin{center}
\begin{tabular}{L{0.16\textwidth} L{0.18\textwidth} L{0.36\textwidth} L{0.12\textwidth}}
\toprule
\textbf{SPR\_ID} & \textbf{Site Type} & \textbf{Name} & \textbf{Land Area} \\
\midrule
416 & Complex & IndotFac - Columbus Lab & 1800 \\
\bottomrule
\end{tabular}
\end{center}

\textbf{Likely reason:} The all-candidate file and scored files come from different
ArcGIS exports. The final candidate layer was exported later from
\code{Cleaned_Solar_Potental_CZ_remove(exclusion)}, while the scored files were exported
from \code{solar_potential_scored_indotfacility} and
\code{solar_potential_scored_interchage}. The Columbus Lab site is very small and may
have been excluded from or never included in the scored subsets.

\textbf{Implementation response:} Include the site in the candidate layer but mark
score fields as missing or \code{not_scored}. Do not silently drop it unless the project
owner confirms that unscored sites should be excluded.

\subsection{Land-Cover Fields Use Different Units}

Fields such as \code{DevOpe_21}, \code{DevLowI_22}, \code{DDevMed_23}, and
\code{CultCr_82} sum to approximately \code{100} in the all-candidate file, but to
approximately \code{1} in the scored files.

\textbf{Likely reason:} The all-candidate workflow calculated land-cover values as
percentages. The scored files appear to use normalized fractions for suitability
scoring. The ArcGIS lineage includes percentage calculations such as
\code{Fld_Perc = (Fld_Area / Shape_Area) * 100}.

\textbf{Implementation response:} Create cleaned fields with explicit unit names,
for example \code{landcover_developed_open_pct}. Convert fractions to percentages
for display and keep raw source fields separately.

\subsection{One Invalid Geometry}

One geometry is invalid in both the all-candidate and ROW scored datasets:

\begin{center}
\begin{tabular}{L{0.18\textwidth} L{0.24\textwidth} L{0.42\textwidth}}
\toprule
\textbf{SPR\_ID} & \textbf{Site Type} & \textbf{Issue} \\
\midrule
371 & Interchange & Ring self-intersection \\
\bottomrule
\end{tabular}
\end{center}

\textbf{Likely reason:} The invalid polygon was probably inherited from the source
GIS layer or introduced by geoprocessing, especially the erase operation used in
the final candidate workflow.

\textbf{Implementation response:} Repair geometries during the cleaning step using
Shapely's geometry validation tools before writing GeoJSON or loading into PostGIS.
Log any repaired records.

\subsection{Layer Naming and Site Labels Are Confusing}

Several ROW/interchange records have display names starting with \code{IndotFac -}.
The \code{Site_typ} field also mixes casing and labels such as \code{Interchange},
\code{unit}, and \code{Complex}.

\textbf{Likely reason:} The ArcGIS metadata shows broad field calculations such as
adding \code{IndotFac -} to site-name fields. This was applied to the prepared GIS
data rather than tailored to final public-facing labels.

\textbf{Implementation response:} Preserve the raw fields but add cleaned display
fields such as \code{site_type_display}, \code{site_name_display}, and
\code{source_layer}. Normalize labels for the UI.

\subsection{Some Suitability Fields Are Layer-Specific}

Facility \code{evp_s} is always \code{0.0}; ROW \code{dem_s} is always \code{0.0}; and
facility \code{dem_s} is limited to a narrow range.

\textbf{Likely reason:} The reference decision-support report explains that different
property types are evaluated using different suitability logic. Some criteria are
not meaningful for every site type.

\textbf{Implementation response:} Do not present every score field as equally
applicable to every layer. The frontend should hide, disable, or explain
non-applicable criteria per layer.

\section{Probable Data Lineage}

The metadata indicates this rough lineage:

\begin{enumerate}
  \item Source GIS data was prepared in ArcGIS Pro from INDOT facilities, rest/welcome areas, interchanges, open land, environmental layers, transmission line proximity, floodplain data, land-cover data, and related suitability factors.
  \item Intermediate candidate sites were cleaned and enriched with derived fields such as land area, flood percentage, forest percentage, transmission distance, and solar mean.
  \item \code{SPR_ID} was calculated from feature IDs during preparation.
  \item Some text labels were modified with prefixes such as \code{IndotFac -}.
  \item Final candidate sites were processed through an erase operation against a road/clear-zone buffer.
  \item The all-candidate shapefile was exported from the later candidate layer.
  \item The scored facility and ROW shapefiles were exported separately from scored subset feature classes.
\end{enumerate}

This explains why the files align closely but are not identical: they are related
exports from the same project, but not a single perfectly synchronized output.

\section{Risks and Mitigations}

\begin{longtable}{L{0.30\textwidth} L{0.30\textwidth} L{0.30\textwidth}}
\toprule
\textbf{Risk} & \textbf{Impact} & \textbf{Mitigation} \\
\midrule
\endhead
Using stale \code{Lat}/\code{Long} fields & Sites would render at the wrong location or fail map display. & Derive coordinates from geometry only. \\
Mixing percentage and fraction units & Misleading land-cover values and filters. & Normalize to explicit cleaned percent fields. \\
Dropping unscored sites silently & Candidate inventory would not match source inventory. & Include unscored status and document missing score fields. \\
Invalid geometry during export & GeoJSON/PostGIS errors or broken map rendering. & Repair and log geometry fixes. \\
Confusing labels in public map & Users may misunderstand facility vs ROW layers. & Add cleaned display labels and preserve raw source values. \\
Layer-specific scores shown globally & Users may compare non-applicable criteria. & Configure score display by site type/layer. \\
\bottomrule
\end{longtable}

\section{Immediate Next Work}

The next implementation step should be a deterministic data-cleaning pipeline that:

\begin{enumerate}
  \item Reads the three shapefiles.
  \item Repairs invalid geometries and logs repaired records.
  \item Converts all geometries to \code{EPSG:4326}.
  \item Calculates centroid latitude and longitude from geometry.
  \item Normalizes land-cover fields to display percentages.
  \item Adds cleaned fields for site type, source layer, site name, score availability, and score applicability.
  \item Exports clean GeoJSON and CSV files for frontend prototyping.
  \item Writes a validation report summarizing row counts, missing scored sites, geometry repairs, and field normalization.
\end{enumerate}

\section{Conclusion}

The initial assessment confirms that the provided GIS data is usable for a web-based
solar suitability map, but it should not be consumed directly by the frontend. The
data is the product of several ArcGIS exports and includes stale coordinate fields,
one missing scored candidate, one invalid geometry, mixed unit conventions, and
labels that need public-facing cleanup. These are manageable issues if the project
starts with a clear cleaning and validation pipeline before API and frontend work.

\end{document}
